% Plan de acción de f1-data-platform a partir del informe de situación (30/09/2026)
% Compilar desde esta carpeta:  latexmk -xelatex -outdir=build plan_accion.tex
\documentclass[11pt,openany]{report}
% Preámbulo del informe de situación de f1-data-platform (compilar con XeLaTeX o latexmk -xelatex)
\usepackage[a4paper,margin=2.4cm,headheight=14pt]{geometry}
\usepackage{fontspec}
\usepackage[spanish,es-noshorthands,es-tabla]{babel}
\setmainfont{Libertinus Serif}
\setsansfont{Fira Sans}[Scale=MatchLowercase]
\setmonofont{Fira Mono}[Scale=0.86]
\usepackage{microtype}
\usepackage[dvipsnames,table]{xcolor}
\usepackage{graphicx}
\usepackage{booktabs,tabularx,longtable,array,multirow}
\usepackage{enumitem}
\setlist{itemsep=2pt,topsep=4pt}
\usepackage{csquotes}
\usepackage{fancyhdr}
\usepackage{titlesec}
\usepackage{caption}
\captionsetup{font=small,labelfont={bf,sf}}
\usepackage{listings}
\usepackage{tcolorbox}
\tcbuselibrary{skins,breakable}
\usepackage{tikz}
\usetikzlibrary{positioning,arrows.meta,shapes.geometric,shapes.multipart,fit,backgrounds,calc,matrix,chains}
\usepackage{pgfplots}
\pgfplotsset{compat=1.18}
\usepackage{amssymb,pifont}
\usepackage{fvextra}
\usepackage[colorlinks,urlcolor=oiblue!75!black,linkcolor=ink,citecolor=ink]{hyperref}
\usepackage{bookmark}

% Colores del proyecto (paleta Okabe-Ito + rojo F1)
\definecolor{f1red}{HTML}{C8102E}
\definecolor{ink}{HTML}{1F2330}
\definecolor{muted}{HTML}{5B6170}
\definecolor{bronze}{HTML}{A0652A}
\definecolor{silver}{HTML}{7C8691}
\definecolor{gold}{HTML}{B8901A}
\definecolor{oiblue}{HTML}{0072B2}
\definecolor{oiorange}{HTML}{E69F00}
\definecolor{oigreen}{HTML}{009E73}
\definecolor{oisky}{HTML}{56B4E9}
\definecolor{oiverm}{HTML}{D55E00}
\definecolor{oipurple}{HTML}{CC79A7}
\definecolor{codebg}{HTML}{F5F6F8}

\titleformat{\chapter}[display]{\sffamily\bfseries\huge\color{ink}}{\color{f1red}\large\MakeUppercase{\chaptertitlename}\ \thechapter}{4pt}{}
\titleformat{\section}{\sffamily\bfseries\Large\color{ink}}{\color{f1red}\thesection}{0.8em}{}
\titleformat{\subsection}{\sffamily\bfseries\large\color{ink}}{\color{muted}\thesubsection}{0.8em}{}
\titleformat{\subsubsection}{\sffamily\bfseries\normalsize\color{ink}}{}{0em}{}
\setcounter{secnumdepth}{2}
\setcounter{tocdepth}{2}

\pagestyle{fancy}
\fancyhf{}
\fancyhead[L]{\sffamily\small\color{muted}\leftmark}
\fancyhead[R]{\sffamily\small\color{muted}f1-data-platform · Informe de situación}
\fancyfoot[C]{\sffamily\small\thepage}
\renewcommand{\headrulewidth}{0.3pt}

\lstset{
  basicstyle=\ttfamily\small,
  backgroundcolor=\color{codebg},
  frame=none,
  breaklines=true,
  columns=fullflexible,
  keepspaces=true,
  showstringspaces=false,
  xleftmargin=6pt,
  framexleftmargin=6pt,
  aboveskip=8pt,
  belowskip=8pt,
  keywordstyle=\color{oiblue}\bfseries,
  commentstyle=\color{muted}\itshape,
  stringstyle=\color{oiverm},
  literate={á}{{\'a}}1 {é}{{\'e}}1 {í}{{\'i}}1 {ó}{{\'o}}1 {ú}{{\'u}}1 {ñ}{{\~n}}1 {Á}{{\'A}}1 {É}{{\'E}}1 {Í}{{\'I}}1 {Ó}{{\'O}}1 {Ú}{{\'U}}1 {Ñ}{{\~N}}1 {ü}{{\"u}}1 {«}{{\guillemotleft}}1 {»}{{\guillemotright}}1 {→}{{$\to$}}1 {×}{{$\times$}}1
}

% Cajas de resumen, recomendación y aviso
\newtcolorbox{resumen}[1][Resumen]{enhanced,breakable,colback=codebg,colframe=ink,boxrule=0.4pt,arc=2pt,
  fonttitle=\sffamily\bfseries,title={#1},left=8pt,right=8pt}
\newtcolorbox{recomendacion}[1][Recomendación]{enhanced,breakable,colback=oigreen!6,colframe=oigreen!70!black,boxrule=0.4pt,arc=2pt,
  fonttitle=\sffamily\bfseries,title={#1},left=8pt,right=8pt}
\newtcolorbox{aviso}[1][Atención]{enhanced,breakable,colback=oiorange!8,colframe=oiorange!80!black,boxrule=0.4pt,arc=2pt,
  fonttitle=\sffamily\bfseries,title={#1},left=8pt,right=8pt}

% Estilos TikZ comunes
\tikzset{
  caja/.style={draw=ink,rounded corners=2pt,align=center,font=\sffamily\small,inner sep=5pt,fill=white},
  capa/.style={draw=none,rounded corners=4pt,inner sep=8pt},
  flecha/.style={-{Stealth[length=2.4mm]},thick,color=ink!80},
  flechad/.style={-{Stealth[length=2.4mm]},thick,dashed,color=muted},
  etiqueta/.style={font=\sffamily\scriptsize,color=muted},
  tabla/.style={rectangle split,rectangle split parts=2,draw=ink,rounded corners=1pt,align=left,font=\sffamily\scriptsize,
    rectangle split part fill={#1,white},inner sep=3pt,text width=3.3cm},
}

\newcommand{\code}[1]{\texttt{\small\color{ink!85!oiblue}#1}}
\newcommand{\urltexto}[1]{\texttt{\footnotesize #1}}
\newcommand{\si}{\textcolor{oigreen!80!black}{\ding{51}}}
\newcommand{\no}{\textcolor{oiverm}{\ding{55}}}
\newcommand{\parcial}{\textcolor{oiorange}{\ding{108}}}
\newcommand{\nada}{\textcolor{muted}{\ding{109}}}

% Bloques de código de las notas (fvextra corta las líneas largas)
\DefineVerbatimEnvironment{codigo}{Verbatim}{breaklines,breakanywhere,fontsize=\footnotesize,
  frame=leftline,framerule=1.2pt,rulecolor=\color{oiblue!50},xleftmargin=8pt,framesep=6pt}
\DefineVerbatimEnvironment{diagramatexto}{Verbatim}{breaklines,breakanywhere,fontsize=\scriptsize,
  frame=single,rulecolor=\color{muted!50},framesep=5pt,label={\sffamily\scriptsize descripción del diagrama}}
\newenvironment{nota}{\begin{tcolorbox}[enhanced,breakable,colback=codebg,colframe=codebg,
  borderline west={2pt}{0pt}{muted},arc=0pt,left=8pt,right=6pt,top=3pt,bottom=3pt,fontupper=\small]}{\end{tcolorbox}}
\newcommand{\fich}[1]{\texttt{\small\color{oiblue!70!black}#1}}
\newcommand{\prio}[1]{\textsf{\small\bfseries #1}}
\setlength{\emergencystretch}{3em}
\counterwithout{figure}{chapter}

\hypersetup{pdftitle={f1-data-platform: plan de acción},pdfauthor={Antonio Sánchez de la Blanca Romero}}
\renewcommand{\thesection}{\arabic{section}}
\counterwithout{section}{chapter}
\fancyhead[L]{\sffamily\small\color{muted}Plan de acción}
\fancyhead[R]{\sffamily\small\color{muted}f1-data-platform · 30/09/2026}

\newcommand{\bloque}[1]{\textsf{\bfseries\color{oiblue!80!black}#1}}

\begin{document}

\begin{titlepage}
\thispagestyle{empty}
\begin{tikzpicture}[remember picture,overlay]
  \fill[ink] (current page.north west) rectangle ([yshift=-5.6cm]current page.north east);
  \fill[f1red] ([yshift=-5.6cm]current page.north west) rectangle ([yshift=-5.8cm]current page.north east);
  \node[anchor=north west,text=white,font=\sffamily\bfseries\fontsize{28}{32}\selectfont]
    at ([xshift=2.4cm,yshift=-2cm]current page.north west) {Plan de acción};
  \node[anchor=north west,text=white!75!ink,font=\sffamily\large]
    at ([xshift=2.4cm,yshift=-3.5cm]current page.north west) {f1-data-platform · alternativas y camino recomendado};
\end{tikzpicture}
\vspace*{5.4cm}

\noindent{\sffamily\color{muted}30 de septiembre de 2026 · basado en el \emph{Informe de situación del
proyecto} (\fich{docs/informe\_situacion/build/informe.pdf})}

\bigskip
\begin{resumen}[En una página]
\begin{itemize}
  \item El informe no encuentra nada roto de gravedad: la plataforma funciona, está desplegada y sus
  pruebas pasan. Lo que hay son \textbf{riesgos operativos baratos de quitar}, \textbf{errores de
  datos concretos}, \textbf{datos disponibles que no se usan} y una web que \textbf{replica bien el
  TFG pero no cuenta aún las carreras}.
  \item Se proponen \textbf{cuatro caminos} alternativos: (A) consolidar y ampliar, (B) escaparate
  visual primero, (C) directo a la fase 7 y (D) portfolio de arquitectura. Todos comparten un
  \textbf{tronco común} de unos 3 días que conviene hacer en cualquier caso.
  \item \textbf{Recomendación: camino A}, en unas 8 semanas de trabajo efectivo, con dos puntos de
  decisión donde se puede derivar hacia B o C. Deja la fase 7 sobre datos sólidos (meteorología,
  compuesto real, neutralizaciones bien marcadas) y la web con las visualizaciones que más valor
  aportan en la defensa.
  \item Necesito de ti \textbf{seis decisiones} (sección~\ref{sec:decisiones}) antes de empezar;
  para el tronco común solo hace falta la primera (licencia).
\end{itemize}
\end{resumen}
\vfill
\end{titlepage}

\tableofcontents
\clearpage

% ======================================================================
\section{Lectura crítica del informe}

Antes de planificar he revisado el informe completo buscando contradicciones entre notas,
dependencias ocultas entre propuestas y cifras que conviene confirmar.

\subsection*{Correcciones y matices}

\begin{itemize}
  \item \textbf{Race trace: desde 1990, no desde 1950.} El catálogo externo lo daba por factible
  desde 1950; en la base no hay tiempos por vuelta ni diferencias antes de 1990 (menos del 1,5\,\%
  por década). Ya está corregido en el informe. El diferencial real es el lap chart desde 1950 y el
  race trace desde 1990.
  \item \textbf{Unit tests de dbt.} Una nota cuenta 4 y otra 2; en el código hay 2 en
  \fich{\_quality\_\_models.yml}. Lo relevante coincide en ambas: ninguno cubre la fusión de
  vueltas ni la numeración, que es la lógica más delicada.
  \item \textbf{Tamaño de la base.} \code{f1.duckdb} pesa 54,5\,MB, no unos 200\,MB. Esto abarata
  cualquier alternativa (meter la base en la imagen Docker, servirla estática, DuckDB-WASM).
  \item \textbf{Cifras de planes gratuitos} (Koyeb, Render, GitHub Pages) proceden en parte de
  agregadores: hay que confirmarlas en la web oficial antes de migrar nada.
  \item \textbf{La API sin caché de página en Vercel} está diagnosticada pero no confirmada: la causa
  probable (rutas dinámicas sin \code{generateStaticParams} o \code{fetch} sin caché) hay que
  verificarla con la salida de \code{next build}.
\end{itemize}

\subsection*{Dependencias que condicionan el orden}

\begin{enumerate}
  \item \textbf{El cambio de grano de \code{fact\_laptimes}} (columna \code{session\_type} para los
  sprints) afecta a todos los consumidores de vueltas: API, race trace, degradación. Conviene
  hacerlo \emph{antes} de construir gráficos o modelos nuevos sobre las vueltas, o todos tendrán que
  revisarse después.
  \item \textbf{La fase 7 depende de tres huecos de datos}: compuesto real C1–C6 desde 2024,
  meteorología y neutralizaciones bien marcadas (la bandera roja cae una vuelta antes en 22
  carreras). Empezar la analítica sin ellos obliga a rehacerla.
  \item \textbf{La arquitectura B} (API estática) reutiliza las consultas de la API. Hacerla antes
  de estabilizar la API (\code{driver\_number}, \code{?session=sprint}) supone exportar dos veces.
  \item \textbf{La navegación de 5 pestañas} decide dónde van los gráficos nuevos. Si se añaden
  gráficos antes de reorganizarla, habrá que moverlos.
  \item \textbf{Integrar Jolpica u OpenF1} requiere antes decidir la licencia (CC BY-NC-SA obliga a
  publicar lo derivado en las mismas condiciones).
\end{enumerate}

\subsection*{Lo que no necesita decisión}

Hay un conjunto de tareas que ningún camino discute y que reducen riesgo inmediatamente. Forman el
\textbf{tronco común} (sección~\ref{sec:tronco}).

% ======================================================================
\section{Bloques de trabajo}

Cada bloque es una unidad que se puede planificar, entregar y verificar por separado. Los caminos
de la sección~\ref{sec:caminos} son combinaciones y órdenes distintos de estos bloques.

\begin{center}\small
\begin{longtable}{@{}>{\raggedright\arraybackslash}p{0.07\linewidth}>{\raggedright\arraybackslash}p{0.46\linewidth}>{\raggedright\arraybackslash}p{0.1\linewidth}>{\raggedright\arraybackslash}p{0.09\linewidth}>{\raggedright\arraybackslash}p{0.18\linewidth}@{}}
\toprule
\textbf{Id} & \textbf{Contenido} & \textbf{Días} & \textbf{Depende} & \textbf{Terminado cuando} \\
\midrule\endhead
\bottomrule\endlastfoot
\bloque{T1} & Higiene: \code{LICENSE} y atribuciones; Dependabot; acciones fijadas a SHA y fuera de Node 20; cabeceras de seguridad; \code{fetch} con timeout y reintento; ruff de pre-commit alineado; región \code{fra1} en Vercel & 1–1,5 & D1 & CI verde sin avisos de Node 20; cabeceras visibles en producción \\
\bloque{T2} & Robustez operativa: releases fechadas e inmutables además de \code{data-latest}; copia del bronze estático; prueba de humo de la API antes de publicar; arranque de la API con la última copia si falla GitHub; \code{/health} con fecha de datos & 1,5–2 & — & Un pipeline con datos rotos no publica; la API arranca sin GitHub \\
\bloque{T3} & Arreglos de datos rápidos: \code{has\_sprint}; bandera roja en la vuelta que contiene la suspensión; tiempo por suma de sectores (570 vueltas); \code{NAN}/\code{NONE} a nulo; \code{HUGEINT}; tabla huérfana; recargar Italia 2018 en FastF1 & 1 & — & QA sin regresiones; São Paulo 2024 v.\,33 marcada \\
\bloque{C1} & Correcciones visibles: \code{driver\_number} en \code{/laps} y \code{/stints}; abreviaturas únicas por carrera; dorsales solo de sesiones de carrera; desempate determinista en \code{int\_lap\_numbering} & 1,5–2 & T3 & Ningún duplicado de grano; Schumacher M./R. distinguibles \\
\bloque{C2} & Pruebas: unit tests de \code{int\_laptimes}, \code{int\_lap\_numbering} y \code{parse\_duration\_ms}; tests de clave y relación en paradas, standings, agregados y telemetría & 2 & C1 & Los tests fallan si se rompe la fusión \\
\bloque{C3} & Errores de lectura de la web: referencia de clasificación explícita; eje continuo de temporadas; constructores sin barras antes de 1958; etiquetas finales sin solaparse; formato español en ejes; paradas con eje ajustado y rótulo «tiempo en pit lane»; violín legible & 2 & — & Revisión visual en escritorio y 375\,px; Playwright + axe en verde \\
\bloque{D1} & FastF1 ampliado: sesiones S/SQ/SS/FP, \code{weather}, \code{messages}; modelos de meteorología y dirección de carrera & 2–3 & T2 & Meteo y mensajes 2018+ en gold \\
\bloque{D2} & Jolpica como segunda fuente (sustituye el volcado Ergast 2022) & 1–2 & D0 (licencia) & 2023–2026 con \code{validation\_status} distinto de \code{single\_source} \\
\bloque{D3} & Eventos históricos SC/VSC/rojas (seed desde RaceData) y nominaciones Pirelli C1–C6 (seed manual) & 3 & — & Banderas rojas 1950+; compuesto real 2024+ \\
\bloque{S1} & Sprints: \code{session\_type} en hechos de vueltas, \code{int\_sprint\_laptimes} con vuelta 1 derivada, \code{?session=sprint} en la API, selector en la web & 4 (mínimo) a 6–8 (con validación) & D1, C2 & Los 29 sprints con lap chart, estrategia y ritmo \\
\bloque{W1} & Navegación: 5 pestañas de carrera (telemetría en Clasificación; Neumáticos + Boxes = Estrategia); pilotos seleccionados en la URL; enlaces cruzados & 3 & C3 & 5 pestañas caben en 375\,px \\
\bloque{W2} & Tanda 1 de gráficos: race trace (1990+); franjas SC/VSC/rojas; matriz de resultados de temporada; récords ampliados; diferencia con el líder y eliminación matemática; KPIs de relato & 6–8 & W1, T3 & Cada gráfico con resumen, tabla y prueba axe \\
\bloque{W3} & Tanda 2: página de circuito; comparador de dos pilotos; parrilla → V1 → llegada; campeonatos con otros sistemas de puntos; detector de undercut & 7–9 & W2 & Rutas nuevas con pruebas e2e \\
\bloque{W4} & Tanda 3 (originales): red de compañeros; ghost race accesible; streamgraph de constructores; ADN del piloto; línea de títulos & 6–10 & W2 & — \\
\bloque{A1} & Arquitectura A: \code{justfile}; sources dbt de todo el bronze con \code{freshness}; subcarpetas en intermediate; disparo diario por calendario y \emph{keepalive} & 2–3 & T2 & \code{just build} reproduce el pipeline local \\
\bloque{A2} & Arquitectura B: \code{export-static}, GitHub Pages, cliente con \code{F1\_DATA\_BASE\_URL} y revalidación bajo demanda & 6–9 & A1, C1, S1 & La web funciona con la API apagada \\
\bloque{F0} & Preparación de la fase 7: vueltas limpias (sin neutralización ni boxes), corrección de combustible, categorías de abandono (seed), telemetría 2018–2023 & 3–4 & D1, D3 & Tabla de vueltas limpias validada \\
\bloque{F1} & Fase 7: degradación, predicción de resultado, estilos de pilotaje & — & F0 & Según su propio plan \\
\end{longtable}
\end{center}

\subsection{Tronco común (en cualquier camino)}\label{sec:tronco}

\textbf{T1 + T2 + T3}, unos 3,5–4,5 días. Son arreglos de bajo riesgo que eliminan los tres puntos
frágiles señalados por el informe (release sobrescrita, publicación sin prueba de humo, falta de
licencia), corrigen los errores de datos más visibles (sprints ocultos, bandera roja) y quitan los
avisos de la CI. Ningún camino gana nada retrasándolos.

% ======================================================================
\section{Caminos alternativos}\label{sec:caminos}

\subsection*{A · Consolidar y ampliar (recomendado)}
\textbf{Secuencia:} tronco → C1 → C2 → D1 → D2 → D3 → S1 → C3 → W1 → W2 → A1 → F0 → (A2 opcional) → F1.\\
\textbf{Duración:} unas 8 semanas hasta F0; A2 añade 1,5–2.\\
\textbf{A favor:} ataca las dependencias en su orden natural (grano de vueltas antes que gráficos y
modelos; datos antes que analítica); cada entrega deja el sistema mejor y verificable; la fase 7
empieza sobre datos que se pueden defender.\\
\textbf{En contra:} la web no cambia de forma visible hasta la semana 4–5.\\
\textbf{Elegirlo si} la prioridad es la solidez del proyecto y la defensa técnica.

\subsection*{B · Escaparate visual primero}
\textbf{Secuencia:} tronco → C3 → W1 → W2 → W3 → C1 → C2 → D1 → S1 → \ldots\\
\textbf{Duración:} unas 4 semanas hasta tener tandas 1 y 2 publicadas.\\
\textbf{A favor:} el cambio se ve enseguida; el race trace y la página de circuito son lo más
vistoso.\\
\textbf{En contra:} los gráficos se construyen sobre un grano de vueltas que cambiará con los
sprints (habrá que revisarlos) y sobre datos con errores aún abiertos (coches compartidos,
abreviaturas); el race trace de 2025–26 hereda los huecos de fuente única.\\
\textbf{Elegirlo si} hay una fecha próxima para enseñar la web (defensa, entrevista, portfolio).

\subsection*{C · Directo a la fase 7}
\textbf{Secuencia:} tronco → D1 → D3 → F0 → F1, dejando la web y los sprints para después.\\
\textbf{Duración:} unas 2,5 semanas hasta empezar la analítica.\\
\textbf{A favor:} la parte más novedosa del proyecto llega antes.\\
\textbf{En contra:} sin C2 la analítica se apoya en una fusión de vueltas sin pruebas; sin S1 el
cambio de grano llega después y obliga a revisar los modelos; la web sigue con sus errores de
lectura.\\
\textbf{Elegirlo si} la fase 7 es el objetivo principal de esta etapa y la web puede esperar.

\subsection*{D · Portfolio de arquitectura}
\textbf{Secuencia:} tronco → A1 → C1 → A2 → (Dagster + R2 + DuckLake experimental) → resto.\\
\textbf{Duración:} 3–4 semanas; 6–8 con la arquitectura C.\\
\textbf{A favor:} elimina el arranque en frío y da un argumento de diseño muy claro; aporta
herramientas valoradas (orquestador, lakehouse).\\
\textbf{En contra:} no mejora los datos ni lo que ve el usuario; la arquitectura C no la exige la
necesidad real y añade proveedores y mantenimiento.\\
\textbf{Elegirlo si} el proyecto se va a usar sobre todo como muestra de ingeniería de datos.

\subsection*{Comparación}

\begin{center}\small
\begin{tabularx}{\linewidth}{@{}>{\raggedright\arraybackslash}p{0.26\linewidth}>{\centering\arraybackslash}X>{\centering\arraybackslash}X>{\centering\arraybackslash}X>{\centering\arraybackslash}X@{}}
\toprule
\textbf{Criterio} & \textbf{A · Consolidar} & \textbf{B · Visual} & \textbf{C · Fase 7} & \textbf{D · Arquitectura} \\
\midrule
Semanas hasta el objetivo & 8 & 4 & 2,5 & 3–4 \\
Cambio visible en la web & semana 4–5 & semana 1–2 & tardío & bajo \\
Calidad de datos al final & alta & media & media & media \\
Retrabajo esperado & bajo & medio (sprints, datos) & medio (grano, tests) & bajo \\
Riesgo técnico & bajo & bajo & medio & medio–alto (con C) \\
Valor para la defensa & muy alto & alto & alto & alto \\
Prepara la fase 7 & sí & no & parcialmente & no \\
\bottomrule
\end{tabularx}
\end{center}

% ======================================================================
\section{Plan recomendado: camino A, semana a semana}

Las semanas son de trabajo efectivo (unos 5 días); si el ritmo es a tiempo parcial, la secuencia es
la misma y solo se estira el calendario.

\begin{figure}[ht]
\centering
\begin{tikzpicture}[x=1.75cm,y=0.55cm,font=\sffamily\scriptsize]
  \foreach \s in {1,...,8} {
    \node[muted] at (\s-0.5,0.7) {S\s};
    \draw[muted!25] (\s-1,0.4) -- (\s-1,-13.4);
  }
  \draw[muted!25] (8,0.4) -- (8,-13.4);
  \newcommand{\barra}[5]{% fila, inicio, fin, color, etiqueta
    \fill[#4,rounded corners=1.5pt] (#2,-#1-0.3) rectangle (#3,-#1+0.3);
    \node[anchor=east,text width=3.3cm,align=right] at (-0.1,-#1) {#5};}
  \barra{1}{0}{0.9}{oiblue}{Tronco: T1 + T2 + T3}
  \barra{2}{0.9}{1.3}{oiblue}{C1 Correcciones visibles}
  \barra{3}{1.3}{1.75}{oiblue}{C2 Pruebas dbt}
  \barra{4}{1.75}{2.3}{oigreen}{D1 FastF1 ampliado}
  \barra{5}{2.3}{2.65}{oigreen}{D2 Jolpica}
  \barra{6}{2.65}{3.25}{oigreen}{D3 Eventos y Pirelli}
  \barra{7}{3.25}{4.4}{oigreen}{S1 Sprints}
  \barra{8}{4.4}{4.8}{oiorange}{C3 Lectura de gráficos}
  \barra{9}{4.8}{5.4}{oiorange}{W1 Navegación}
  \barra{10}{5.4}{6.9}{oiorange}{W2 Tanda 1 de gráficos}
  \barra{11}{6.9}{7.45}{oipurple}{A1 Arquitectura A}
  \barra{12}{7.45}{8}{oipurple}{F0 Preparar fase 7}
  \foreach \x/\t in {1.75/Hito 1,3.25/Hito 2,6.9/Hito 3,8/Hito 4} {
    \draw[f1red,thick] (\x,0.4) -- (\x,-13);
    \node[f1red,font=\sffamily\bfseries\scriptsize,fill=white,inner sep=1pt] at (\x,-13.2) {\t};
  }
\end{tikzpicture}
\caption{Calendario orientativo del camino A. Azul: calidad; verde: datos; naranja: web; morado:
arquitectura y preparación de la fase 7. A2 (API estática) y las tandas W3 y W4 quedan detrás, según
lo que se decida en los hitos.}
\end{figure}

\begin{description}[style=nextline,leftmargin=1.2em,font=\sffamily\bfseries]
  \item[Semana 1 · Tronco común]
  T1, T2 y T3. Al acabar, el pipeline no publica una base rota, las releases quedan fechadas, la API
  arranca aunque GitHub falle, la CI no avisa de Node 20 y los 12 sprints de 2021–2023 aparecen en la
  web.
  \item[Semana 2 · Corrección y pruebas → Hito 1: «plataforma fiable»]
  C1 y C2. Punto de decisión: si surge una fecha para enseñar la web, derivar al camino B desde
  aquí (C3 → W1 → W2) y retomar los datos después.
  \item[Semanas 2–3 · Datos nuevos → Hito 2: «datos completos 2018+»]
  D1, D2 y D3. Al acabar: meteorología y dirección de carrera desde 2018, segunda fuente en
  2023–2026, banderas rojas históricas y compuesto real. Punto de decisión: si la prioridad pasa a
  ser la fase 7, saltar a F0 desde aquí (camino C con base sólida).
  \item[Semanas 4–5 · Sprints y web]
  S1 (con el cambio de grano de vueltas hecho antes de cualquier gráfico nuevo), C3 y W1.
  \item[Semanas 5–7 · Tanda 1 de gráficos → Hito 3: «la web cuenta las carreras»]
  Race trace, franjas de neutralización, matriz de resultados, récords ampliados, eliminación
  matemática y KPIs de relato.
  \item[Semana 8 · Arquitectura A y preparación de la fase 7 → Hito 4]
  \code{justfile}, sources completos, disparo por calendario; vueltas limpias y variables para la
  degradación. A partir de aquí: fase 7, o A2 si el arranque en frío sigue molestando.
\end{description}

\section{Variantes de alcance}

\begin{center}\small
\begin{tabularx}{\linewidth}{@{}l>{\raggedright\arraybackslash}X>{\raggedright\arraybackslash}p{0.17\linewidth}@{}}
\toprule
\textbf{Alcance} & \textbf{Qué incluye} & \textbf{Duración} \\
\midrule
Mínimo & Tronco común + C1 + C3 + sprints mínimos (S1 solo FastF1) & 2–2,5 semanas \\
Estándar (recomendado) & Camino A hasta el hito 4 & 7–8 semanas \\
Ampliado & Estándar + A2 (API estática) + W3 (circuito, comparador) & 11–13 semanas \\
\bottomrule
\end{tabularx}
\end{center}

% ======================================================================
\section{Decisiones que necesito de ti}\label{sec:decisiones}

\begin{enumerate}
  \item \textbf{Licencia} (D0, necesaria para el tronco): propuesta MIT para el código y CC BY 4.0
  para los datos derivados, con atribución a F1DB, FastF1 y los datos © F1/FIA.
  \item \textbf{¿Aceptas la cláusula no comercial (CC BY-NC-SA)} que implicaría integrar Jolpica y
  OpenF1? Si no, D2 se sustituye por validación interna sin publicar esos datos.
  \item \textbf{Permiso para descargar volcados de terceros}: Jolpica (unos 14,5\,MB) y, solo para
  validación interna, \code{f1resultsdatabase} (74\,MB).
  \item \textbf{Camino}: A (recomendado), B, C o D, y alcance (mínimo, estándar o ampliado).
  \item \textbf{Cambio de región de Vercel a Frankfurt} (\code{fra1}): se hace desde el panel o con
  \code{vercel.json}; afecta al despliegue en producción.
  \item \textbf{Navegación de 5 pestañas} y secciones nuevas (Circuitos, Comparar, Historia): ¿de
  acuerdo con la propuesta o prefieres mantener la estructura del TFG?
\end{enumerate}

\noindent Además, quedan en tu mano: publicar la propuesta de correcciones a F1DB, activar la
documentación de dbt en GitHub Pages y, si se elige la variante ampliada, confirmar los límites de
los planes gratuitos en las webs oficiales.

\section{Riesgos y cómo mitigarlos}

\begin{center}\small
\begin{tabularx}{\linewidth}{@{}>{\raggedright\arraybackslash}p{0.3\linewidth}>{\raggedright\arraybackslash}X@{}}
\toprule
\textbf{Riesgo} & \textbf{Mitigación} \\
\midrule
El cambio de grano de \code{fact\_laptimes} rompe consultas existentes & \code{session\_type} con valor por defecto \code{race}; filtro explícito en todas las consultas de la API; test de grano; hacerlo antes de W2 \\
El backfill de FastF1 choca con el límite de 500 peticiones por hora & Cargas por temporada repartidas en varias ejecuciones del pipeline; caché persistente \\
Jolpica u OpenF1 cambian o desaparecen & Guardarlos en bronze versionado (releases fechadas); son fuentes de contraste, no primarias \\
La reclasificación de banderas rojas altera datos ya validados & QA antes y después; lista de las 22 carreras afectadas en el informe de calidad \\
Deriva de alcance en los gráficos & Tandas cerradas (W2, W3, W4) con criterio de terminado; nada de W3 hasta cerrar W2 \\
Tiempo de desarrollo menor del previsto & El camino A se puede cortar en cualquier hito dejando el sistema coherente \\
\bottomrule
\end{tabularx}
\end{center}

\section{Qué no hacer, por ahora}

\begin{itemize}
  \item Exportación estática completa de Next.js (unos 16\,400 HTML; rompe los formularios sin JS).
  \item Migrar a Fly.io o Cloudflare Workers (sin plan gratuito o sin sitio para DuckDB).
  \item Nx o Turborepo (sobredimensionados).
  \item Buscar tiempos por vuelta anteriores a 1990: no hay fuente abierta; se documenta como
  limitación.
  \item Scraping de formula1db.com o statsf1.com: sus condiciones lo impiden.
  \item Arquitectura C (Dagster + R2 + DuckLake) antes de terminar el camino elegido.
\end{itemize}

\end{document}
